\documentclass[conference]{IEEEtran}
\IEEEoverridecommandlockouts
% The preceding line is only needed to identify funding in the first footnote. If that is unneeded, please comment it out.
%Template version as of 6/27/2024

\usepackage{cite}
\usepackage{amsmath,amssymb,amsfonts}
\usepackage{algorithmic}
\usepackage{graphicx}
\usepackage{textcomp}
\usepackage{xcolor}
\usepackage{array}
\usepackage{tikz}
\usetikzlibrary{positioning,arrows.meta}
\def\BibTeX{{\rm B\kern-.05em{\sc i\kern-.025em b}\kern-.08em
    T\kern-.1667em\lower.7ex\hbox{E}\kern-.125emX}}
\newcommand{\pll}[1]{\texttt{#1}}
\begin{document}

\title{An Ontology Based Approach to Scene Generation\\
% \thanks{Identify applicable funding agency here. If none, delete this.}
}

\author{\IEEEauthorblockN{C\"uneyt Yavuz}
\IEEEauthorblockA{Semantic Space \\
Istanbul, Turkey \\
cuneyt.tyler@semanticspace.io}
}

\maketitle

\begin{abstract}
Scene generation for game development with large language models (LLMs) has attracted
growing attention. Recent methods use fine-tuning and reinforcement learning on
generation tasks, or extensive prompting over well-defined modules that describe the
scene to be generated. These works have improved generation quality considerably, but
they lack a structured schema for the vocabulary of the environment and of the objects
placed in it. We propose GBO, the GenBuilding Ontology, which provides the vocabulary and
definitions for building a 2D scene: build types, place modules with typed quantity and
quality fields, asset roles, spatial relations with score functions, and design rules,
each fact with its source. Our generation procedure is built on this ontology. It turns a
request into a sequence of typed, validated documents: a Brief confirmed by the user, a
Spec plan tree, and a program in the Place Layout Language (PLL), a declarative
constraint language that extends GLDL with hierarchy, n-ary group relations, networks,
linear features, subdivision, scatter and terrain. Seeded solvers turn the program into
cells and entities, and each level is verified and repaired locally before the next is
built inside it. An asset manager fills every slot from the project, a curated library,
the web or image generation, according to a fixed decision procedure. We evaluate the
method by ablating components and ranking the generated places by human judgment.
\end{abstract}

\begin{IEEEkeywords}
2D scene generation, game development, ontology, large language models, procedural content generation
\end{IEEEkeywords}

\section{Introduction}
Computational creation of 2D and 3D assets and scenes has been an active area since the
early 2000s. Early methods rely on rule-based mechanisms such as grammars, L-systems and
tile constraint solvers to generate artificial content \cite{parish2001,shaker2016pcg}.
These methods need hand-written rules for every kind of place, and human intervention
and validation of their output, which makes them time-consuming for large tasks. Since
the arrival of LLMs, methods have explored prompt engineering \cite{todd2023,sudhakaran2023},
dividing the building task into separate modules \cite{holodeck}, constraining the
generation to narrow down the work \cite{scenecraft}, and using an intermediate
language that a solver evaluates \cite{gldl}.

These works share one principle: the language model should decide \emph{what} a place
contains and how its parts relate, while numerical solvers decide \emph{where} things
go. LLMs are poor at coordinates and good at relations. What they lack is a shared,
explicit account of the domain. Each system defines its own vocabulary of objects,
relations and parameters inside prompts or learned programs, so the meaning of
``a large village'', ``a garden behind each house'' or ``common flowers'' is implicit,
untyped and not checkable.

Our work proposes a well-defined methodology, a prompting schema, GBO (the GenBuilding
Ontology), PLL (the Place Layout Language) and an asset management procedure for
building a game environment. The methodology is a fixed sequence of steps,
\emph{interpret, confirm, plan, compile, solve, verify}, each of which writes a typed,
constrained document. The prompting schema is the structured output of the LLM during
consecutive steps of the building procedure, with a free-text reasoning phase followed by
a strict answer that is validated against the schema. GBO defines the vocabulary and
properties for the building task; it is consulted at every step of the generation
pipeline. PLL is based on GLDL \cite{gldl} but extends it with a hierarchy of regions,
nested places, n-ary group relations, road and river networks, area subdivision, linear
features on cell edges, scatter and terrain, and provenance links back to the plan.
Asset management takes into account the design principles of the scene to be generated
and makes use of existing web sources as well as generation of assets with AI.

Our contributions are:
\begin{itemize}
\item An explicit place ontology that both \emph{grounds the LLM's planning} (only the
  relevant subgraph is given to each call) and \emph{defines the constraints that are
  compiled and checked}. To our knowledge, no prior system uses one shared ontology for
  both purposes.
\item A planning scheme in which the LLM decides only the upper levels of a place, and
  ontology-driven expansion rules fill the lower levels, so a town needs tens of model
  calls rather than thousands.
\item PLL, a declarative layout language for 2D grid worlds with a fixed, deterministic
  solving order, relation score functions and in-document checks.
\item An asset procedure that matches catalogued sprites to ontology roles with hard
  filters (perspective, footprint, licence) and a fit score, and generates or fetches the
  rest.
\end{itemize}
We evaluated our method by removing certain components from our algorithm and
generating places with the same query. We used human evaluation to rank the generated
scenes.

\section{Related Work}

\subsection{Procedural Generation of Places}
Procedural content generation has developed strong algorithms for the geometric layer
of places \cite{shaker2016pcg,togelius2011}. Parish and M\"uller grow street networks
with extended L-systems \cite{parish2001}. Emilien et al.\ grow villages on arbitrary
terrain from \emph{interest functions}: weighted sums of criteria such as slope, water,
road access and sociability, in which any criterion of $-1$ vetoes a site; roads are
least-cost paths on which reusing an existing road is nearly free \cite{emilien2012}.
Vanegas et al.\ subdivide blocks into parcels with straight-skeleton and
oriented-bounding-box splits that guarantee street access, and map parcels across edits
so that user customisations survive \cite{vanegas2012}. Wave Function Collapse treats
tile synthesis as constraint solving \cite{karth2017}; its hierarchical semantic variant
collapses a weighted ``consists-of'' DAG of meta-tiles from the root down and stores the
collapse paths, so all instances of one meta-tile can be selected and regenerated
\cite{alaka2023}. Dormans generates action-adventure levels as two structures, a mission
graph of what the player does and a space graph of where it happens \cite{dormans2010}.
These algorithms produce good geometry, but each needs its parameters set by hand for
every kind of place.

\subsection{LLM-Driven Scene and Level Generation}
Text-to-level work first prompted or fine-tuned language models to emit tile maps
directly \cite{todd2023,sudhakaran2023}. Later systems separate semantics from geometry.
LayoutGPT asks the model for layouts in a CSS-like format \cite{layoutgpt}. Holodeck runs
a fixed chain of modules (floor and walls, doors and windows, object selection, layout),
one LLM query per module; the model writes spatial \emph{relations}, not coordinates, and
a depth-first grid solver places large anchor objects first \cite{holodeck}. Infinigen
Indoors states constraints over abstract object classes, including accessibility
\cite{infinigen}. SceneCraft represents a scene as a bipartite graph of assets and
relation nodes, each relation linking any subset of assets and having a score function
in $[0,1]$; removing the graph was its most damaging ablation \cite{scenecraft}.
3D-GPT drives documented procedural generators whose outputs remain editable parameters
\cite{3dgpt}, and WorldCraft splits a request into sub-spaces handled by specialist
agents, keeping a manual of verified examples \cite{worldcraft}. CityGenAgent generates
cities as block and building programs, trained with supervised fine-tuning and
reinforcement learning, and edits them by language after generation
\cite{citygenagent}.

Closest to our language, GLDL (\emph{Constraint Is All You Need}) turns a narrative into
definition sentences (predefined facilities, new facilities, points) and constraint
sentences (around a position, by a wall, near, far, can-see, alignment) and solves them
as a facility layout problem with penalties for overlap, a cluster term and a sparsity
term; agent simulation still found unreachable keys in 9.5\% of levels \cite{gldl}.
The same authors' database-driven framework moves the LLM's knowledge into
offline databases of facilities, rooms and mechanics reviewed by experts, so that
assembly itself needs no model calls, and repairs levels in two phases
\cite{xu2025database}. For 2D games, Narrative-to-Scene builds layered tile scenes from
stories and reports that without conflict resolution only 72\% of relations held
\cite{chen2025narrative}, and GameTileNet labels low-resolution tiles with affordances
and tile types \cite{gametilenet}.

\subsection{Semantics and Ontologies for Game Worlds}
Semantic game worlds attach classes, attributes and placement rules to objects.
Tutenel et al.\ use such a semantic library to guide layout solving \cite{tutenel2009},
and SketchaWorld declares a virtual world through a semantic model that procedural
methods then realise \cite{smelik2011}. These libraries are tied to their solvers and
were not designed to ground a language model. On the LLM side, the systems above use
ad-hoc prompts or learned programs. GBO places one ontology between the two: published
in OWL~2 \cite{owl2}, SKOS \cite{skos} and SHACL \cite{shacl}, it is both the vocabulary
the LLM plans with and the source of the constraints the compiler emits and the
verifier checks.

\subsection{Clarification and Structured Output}
Under-specified requests are common in generation. ClarEval classifies ambiguity into
missing goals, missing premises and ambiguous terminology, and finds that agents usually
guess silently \cite{clareval}. Active task disambiguation samples several
interpretations and asks the questions whose answers split them most evenly
\cite{kobalczyk2025}. On output format, Tam et al.\ show that strict format restrictions
during reasoning degrade it \cite{tam2024}. Our confirmation and prompting schemes build
on both findings.

\section{Method}

\subsection{Methodology}

Every build runs the same procedure (Fig.~\ref{fig:pipeline}). Each step reads only
documents written before it and writes one typed document, kept with the place, so any
step can be re-run alone and every built entity can be traced back to the plan node that
made it.

\begin{figure}[t]
\centering
\begin{tikzpicture}[
  step/.style={draw, rectangle, minimum width=1.75cm, minimum height=0.85cm,
               align=center, font=\footnotesize},
  arr/.style={-{Stealth[length=1.6mm]}},
  doc/.style={font=\tiny\itshape}]
\node[step] (int) at (0,0)    {Interpret\\[-1pt]{\tiny Brief drafts}};
\node[step] (con) at (2.25,0) {Confirm\\[-1pt]{\tiny Brief}};
\node[step] (ast) at (4.5,0)  {Assets\\[-1pt]{\tiny asset plan}};
\node[step] (pln) at (6.75,0) {Plan\\[-1pt]{\tiny Spec tree}};
\node[step] (cmp) at (6.75,-1.5) {Compile\\[-1pt]{\tiny PLL program}};
\node[step] (slv) at (4.5,-1.5)  {Solve \& build\\[-1pt]{\tiny scene}};
\node[step] (ver) at (2.25,-1.5) {Verify\\[-1pt]{\tiny report}};
\node[font=\footnotesize, align=center] (out) at (0,-1.5) {place +\\report};
\draw[arr] ([xshift=0.2cm, yshift=0.55cm]int.north west) node[above, font=\footnotesize]{query}
  -- ([xshift=0.2cm]int.north west);
\draw[arr] (int) -- (con);
\draw[arr] (con) -- (ast);
\draw[arr] (ast) -- (pln);
\draw[arr] (pln) -- (cmp);
\draw[arr] (cmp) -- (slv);
\draw[arr] (slv) -- (ver);
\draw[arr] (ver) -- (out);
\draw[arr, dashed] (con.north) to[bend right=35] node[above, font=\tiny]{questions} (int.north);
\draw[arr, dashed] (ver.south) to[bend right=25] node[below, font=\tiny]{local repair} (cmp.south);
\node[draw, rectangle, minimum width=8.5cm, minimum height=0.6cm, font=\footnotesize,
      fill=gray!12] (gbo) at (3.375,-3.05)
      {GBO: build types, modules, fields, roles, relations, rules};
\foreach \n in {ver,slv,cmp} \draw[dotted] (\n.south) -- (\n.south |- gbo.north);
\end{tikzpicture}
\caption{The building procedure. Each step writes a typed document; every step queries
the GBO ontology. Verification failures are repaired locally at the node that caused
them.}
\label{fig:pipeline}
\end{figure}

\paragraph{Interpret} The LLM reads the query and drafts several candidate Briefs. A
Brief holds the build type, the scale as a class \emph{and} numbers, the style (era,
culture, biome, condition), the components (the top-level modules of the build type, each
with a count or density and an asset policy), the layout, terrain features and gameplay
hooks. Each field carries a confidence and its source (user, AI, default or derived).

\paragraph{Confirm} Following the clarification results above, the user confirms the
Brief in a window generated from its schema, in at most two batched rounds. Fields on
which the drafts disagree, or with low confidence, are highlighted as questions, and
fields on which they agree are pre-filled. Vague size words are always shown with their
numbers: the scale table of a settlement maps \emph{village} to 15--60 buildings,
\emph{large village} to 60--150 and \emph{town} to 150--500, and if the user asks for 200
buildings in a ``village'', the class follows the numbers.

\paragraph{Assets} From the confirmed Brief, the asset manager collects every asset the
build will need and decides its source (Section~\ref{sec:assets}).

\paragraph{Plan} The plan is a \emph{Spec}: a tree of module instances, each carrying the
module's quantity fields (counts, areas, sizes, shares, densities with numeric meaning)
and quality fields (values of closed vocabularies). The LLM decides the upper levels in
the Brief's terms, for example the settlement, its wards and per-ward totals and
qualities. Lower levels are filled by the modules' \emph{expansion rules}: defaults and
weighted distributions from the ontology, keyed on the build type, scale and the
parent's qualities, in the spirit of the offline databases of \cite{xu2025database} and
Emilien's building encyclopedia \cite{emilien2012}. A Spec node may be a template:
\texttt{repeat} gives $n$ instances and \texttt{share} assigns a child to a fraction of the
parent's instances, so 28 house plots are one node until the solver must tell them
apart. The LLM is called for a lower node only where its quality fields need judgment,
such as a named landmark.

\paragraph{Compile} Each module has a compile template written in PLL with placeholders
for the node's values. Compiling instantiates the templates top-down, emits the module's
default relations, and adds the relations the LLM proposed for the upper levels. Every
statement records the Spec node it came from.

\paragraph{Solve and build} Seeded solvers evaluate the PLL program region by region
(Section~\ref{sec:pll}) and write a solved layout, which the engine applies as bulk,
undoable commands on chunked tile layers.

\paragraph{Verify and repair} Each level is verified before the next is built inside it.
Hard rules are checked exactly: footprints inside parents, no overlaps on a layer,
allowed terrain, reachability of every door and gate, connected roads, continuous linear
features, perspective compatibility and allowed licences. Conformity with the Spec is
measured against tolerances: counts within $\pm5\%$ (exact below 10), areas within
$\pm10\%$, densities within one level. Finally, the part is rendered headless and a
vision model answers the module's checklist, since the render--critique loop gave
SceneCraft its largest gain \cite{scenecraft}. A failure returns to the stage that made
the part, with the violated rule as the reason, and is repaired at that node only, in
the two-phase order of \cite{xu2025database}: move adaptable results, remove adaptable
results that block a hard rule, re-solve the subtree with a new seed, and finally
re-plan the node. Every approved module build is stored as a verified example for later
prompts, as in WorldCraft's manual \cite{worldcraft}.

\paragraph{Editing after generation} Every built entity keeps a link to its place, Spec
node, PLL statement and instance. Changing a field of a node changes the Spec and
recompiles only that subtree, with siblings and parents held fixed. A hand edit inside a
node locks it, and its entities become fixed statements on the next regeneration. When a
parent is regenerated, old children are matched to new ones by id and position, as in
parcel persistence \cite{vanegas2012}, so locks and edits survive.

\paragraph{Build types} Plan and Compile proceed through stages that are fixed per build
type, in the manner of Holodeck's modules \cite{holodeck}. A settlement runs \emph{site
and terrain $\rightarrow$ networks (main roads, river, walls and gates) $\rightarrow$
wards $\rightarrow$ streets and plots $\rightarrow$ parcel contents $\rightarrow$
dressing}; an interior runs \emph{room programme $\rightarrow$ floor plan $\rightarrow$
doors and windows $\rightarrow$ furniture groups $\rightarrow$ small objects $\rightarrow$
decor and lights}; a dungeon runs \emph{mission graph $\rightarrow$ space graph
$\rightarrow$ geometry $\rightarrow$ room contents $\rightarrow$ encounters}
\cite{dormans2010}. A region plans its settlements, castles and dungeons as nested places,
each with its own documents, solved inside the boundary and entry anchors its parent
found for it. The style is fixed in the Brief before any stage, so every module and
asset choice agrees.

\subsection{Prompting Schema}

Every LLM call in the procedure has the same shape:
\begin{equation*}
\text{inputs} \rightarrow \text{think} \rightarrow \text{answer} \rightarrow
\text{validate} \rightarrow \text{(repair once)}.
\end{equation*}

\paragraph{Inputs} The builder assembles, per call: the step's task description and its
output schema; a one-shot example, following Holodeck's prompt structure \cite{holodeck};
the \emph{relevant ontology subgraph} (the node's module, its allowed children with
weights, the rules and relations that apply to it, and the matching concepts and roles);
the documents written so far (the Brief and the parent Spec node); and verified examples
of similar nodes. The whole ontology is never placed in a prompt.

\paragraph{Think, then answer} Planning a place is multi-step arithmetic and rule
checking: child areas must fit their parent, shares must match counts, and design rules
interact. Because strict formats degrade reasoning \cite{tam2024}, the model first
reasons in free text: what the request implies, which rules from the subgraph apply,
the numbers worked through, and the alternatives considered. Only then does it write the
step's document in its strict schema. The task description ends with a self-check list
that the model works through while thinking; for the Plan step it asks whether the
children's areas fit within 85\% of the parent, whether counts sum to the Brief's numbers,
whether every enabled component is placed, and whether each zone follows the ontology's
site preferences. The answer includes a short \emph{rationale} for each decided part,
which the editor shows when the user asks why a part was built as it was, together with
the ontology facts the call used. The raw reasoning is logged but never shown.

The reasoning is produced in one of three modes, chosen per provider: \emph{native}, when
a reasoning model returns its thinking separately; \emph{two-pass}, the default, in which a
first call reasons without a format and a second call writes the answer given the inputs
and the reasoning; and \emph{one-pass}, in which the schema's first field is a
\texttt{thinking} string. Steps that resolve ambiguity or plan (Interpret, the upper Plan
levels, upper-level relations in Compile, visual critique and re-planning) think;
expansion of lower levels makes no model call at all.

\paragraph{Validate and repair} Each answer is validated against the schema, the tree
structure and the aggregate constraints. A failed answer gets one retry with the exact
violations listed; if that also fails, the step reports the problem rather than
guessing.

\paragraph{Typed fields} The schemas restrict what an answer may express. Every field has
a unit, a range or a vocabulary, a default and the bookkeeping flags \texttt{source},
\texttt{confidence} and \texttt{locked}. Quantity types are \texttt{count}, \texttt{area},
\texttt{size}, \texttt{length}, \texttt{width}, \texttt{spacing}, \texttt{ratio},
\texttt{share} and \texttt{density}. A density is one of five words (\emph{none, rare,
some, common, abundant}) that each field maps to numbers, for instance $0, 0.05, 0.15,
0.3, 0.5$ flowers per free cell, so the word is never ambiguous. Quality types are
\texttt{enum}, \texttt{enumSet}, \texttt{mix} (weighted vocabulary values summing to 1)
and \texttt{bool}, plus \texttt{tags} and free \texttt{description}. The central rule of
the schema is that \emph{decisions that change the layout are made only through
quantity, enum, enumSet, mix or bool fields}, which are closed and checkable; tags and
descriptions are used only for asset retrieval, generation and the critic.

Aggregate constraints across the tree are written in a small expression language with
paths (\texttt{self}, \texttt{parent}, \texttt{children(m)}, \texttt{descendants(class: C)}),
aggregates, arithmetic, comparison and logic, for example
\begin{equation*}
\texttt{sum(children().area) <= 0.85 * self.area}
\end{equation*}
for a settlement, which leaves room for streets. During planning a field may still hold a
range; a constraint over ranges is satisfiable when some choice within the ranges meets
it, which is checked with interval arithmetic. The compiler then picks seeded values that
meet every constraint of the node and records them. Constraints that depend on where
things land, such as footprint density or reachability, are not plan-time constraints;
they become PLL checks.

\subsection{GBO: GenBuilding Ontology}

Module definitions describe each part on its own, but the builder also needs facts that
cross modules: a water mill \emph{requires} a river; a tannery is \emph{typically
downstream of} houses; \emph{smithy}, \emph{forge} and \emph{blacksmith} name the same
concept; a well is \emph{a kind of} water source; an asset showing a rose \emph{can fill}
the flower role; patricians \emph{prefer} high ground. GBO holds all of them in one
queryable graph.

\paragraph{Classes} The terminology contains nine classes. \texttt{BuildType} is a kind
of place (settlement, interior, dungeon, region) with its stages, Brief schema and scale
table of \texttt{ScaleClass} entries. \texttt{Module} is a part at a \emph{level}, from
outer to inner \emph{site, zone, parcel, feature, dressing}; for example
\texttt{Garden}~$\sqsubseteq$~\texttt{Feature} and
\texttt{ResidentialWard}~$\sqsubseteq$~\texttt{Zone}. Modules are further grouped into
module classes such as \texttt{BuildingParcel} or \texttt{PlantedArea}, so constraints can
count ``all parcels whose primary object is a building'' without listing module names.
\texttt{Field} is a quantity or quality field of a module, \texttt{Concept} a vocabulary
value (\texttt{Rose}~$\sqsubseteq$~\texttt{Flower}), \texttt{Role} an asset role a slot
needs (\texttt{FencePieceSet}, \texttt{HouseBody}), \texttt{Asset} a catalogued asset of
the project, \texttt{Relation} a spatial relation type with its score function, and
\texttt{Rule} a design rule with its source.

\paragraph{Properties} Table~\ref{tab:gbo} lists the main properties. The
\texttt{consistsOf} edges with weights and cardinalities form the weighted DAG of
hierarchical WFC \cite{alaka2023}: densities are edge weights, and constraints written
at a parent apply to its subtree. \texttt{prefersSite} carries Emilien's interest
criteria \cite{emilien2012}. \texttt{canFill} links the project's catalogued assets to
roles. Every fact carries a \texttt{source}: a paper, a design rule, a user or a
verified build.

\begin{table}[t]
\caption{Main GBO properties}
\label{tab:gbo}
\centering
\footnotesize
\begin{tabular}{@{}>{\ttfamily}p{1.75cm} p{1.75cm} p{4.1cm}@{}}
\hline
\textrm{\textbf{Property}} & \textbf{Between} & \textbf{Carries} \\
\hline
subClassOf & classes & taxonomy (Garden $\sqsubseteq$ Feature) \\
consistsOf & module $\rightarrow$ module & weight (typical share), min/max cardinality, condition \\
requires, excludes & module, feature & WaterMill requires River \\
typically\-Related & module $\times$ module & a relation with a weight (Tannery downstream-of houses) \\
prefersSite & module $\rightarrow$ criterion & signed weight (Church: high ground $+0.8$) \\
hasField & module $\rightarrow$ field & type, range, default, density map \\
needsRole & module $\rightarrow$ role & essential or adaptable; joins \\
canFill & asset $\rightarrow$ role & confidence, perspective, footprint \\
label, altLabel & any $\rightarrow$ text & names and synonyms, for interpretation \\
broader, narrower & concept $\leftrightarrow$ concept & SKOS hierarchy of values \\
source & fact $\rightarrow$ source & paper, rule, user or verified build \\
\hline
\end{tabular}
\end{table}

\paragraph{Modules} A place module is data: its quantity and quality fields, its asset
\emph{slots}, its allowed children with typical weights, its aggregate constraints, its
default relations and its compile template. A slot binds to an ontology role and names
the quality field whose values say which concepts the art must show; bindings may be
conditional. The garden module, for example, binds its \texttt{plant} slot to the
\texttt{crop} role when its \texttt{kind} is \emph{vegetable} and to \texttt{flower}
with the concepts of its \texttt{flower-mix} when it is \emph{flower}; its \texttt{bed}
slot is adaptable, with tilled ground as the fallback. Each slot is marked
\emph{essential} or \emph{adaptable}, which tells repair what it may move or drop. A
module lists its children; its possible parents are derived, never written twice.
The first essential slot of a parcel or feature (a house, a well, a dining table) is its
\emph{primary} object, which children refer to as \texttt{parent.primary}.

\paragraph{How each step uses it} The interpreter maps the query's words to concepts
through labels, synonyms and \texttt{broader} links, and chooses the build type. Confirm
takes from it the allowed components, the fields to ask and the meaning of size words.
Plan draws allowed children, weights, requirements, exclusions and site preferences.
Compile takes the geometric meaning of each relation. Asset retrieval asks which assets
fill a role in the grid's perspective. Verification asks which rules apply to a node and
its subtree, inherited through the class hierarchy. Explanations cite the facts a
decision used.

\paragraph{Representation} Inside the engine, GBO is a typed property graph loaded from
JSON files with stable identifiers; module files are read as \texttt{Module} classes, so
there is one source of truth. It is published in three Turtle files under a persistent
namespace: the core vocabulary in OWL~2, written by hand; the library in OWL~2 and SKOS,
with modules punned as classes under their level, child cardinalities as OWL
restrictions on \texttt{hasPart}, and vocabularies and roles as SKOS concept schemes; and
SHACL shapes for the library itself and for built places. The library and shapes are
generated from the JSON, and a test fails when the published files fall behind. The
aggregate arithmetic and conditional bindings go beyond SHACL Core, so they are carried
as expression strings and evaluated by the engine.

The current settlement library holds 72 modules with 202 fields, 99 slots and 120 child
rules; 77 asset roles; 20 vocabularies with 107 concepts; 26 relation types; and 14
design rules, each with its source (for example ``the market sits at the centre, on the
main road'' from Watabou's city generator \cite{watabou}, and ``patricians live on high
ground'' from \cite{emilien2012}; others follow Lynch's elements of the city image
\cite{lynch1960}).

\paragraph{Growth} The core is authored by hand from the literature. New modules or
concepts (a lighthouse, say) can be proposed by the AI as small subgraphs, validated like
any document and added after review, which is the offline authoring of
\cite{xu2025database}. Approved builds adjust the \texttt{consistsOf} and
\texttt{typicallyRelated} weights; such facts are marked \texttt{source: learned} and kept
separable from authored ones. Projects may add their own concepts, modules and roles
without changing the engine's core.

\subsection{Place Layout Language (PLL)}
\label{sec:pll}

PLL is the intermediate language between the Spec (\emph{what and how much}) and the
scene (\emph{cells and entities}). A PLL document says where things go and by which
rules. Its design goals are: declarative statements; static checkability before solving
(every reference resolves, every value has a type and unit, and children too large for
their parent are caught); determinism (the same document and seed give the same place);
traceability to Spec nodes; hierarchy and locality (a subtree can be re-solved alone with
everything outside it fixed); native 2D grid semantics; and a small core extended
through the ontology.

\paragraph{Space} Positions are integer cells; an \emph{edge} is the boundary between two
adjacent cells, on which fences and walls lie. Directions are grid directions, also on
isometric grids. Objects have a footprint taken from the catalog (the opaque base of
their art, not the image size), a facing and a layer (\emph{ground, decal, prop,
building, overhead}); footprints may overlap only across layers. Distance is the
smallest Chebyshev distance between the cells of two results.

\paragraph{Statements} A document is a flat, unordered list of nineteen statement kinds
(Table~\ref{tab:pll}). Each has an id and the Spec \texttt{node} it came from, and may be
marked essential. A reference names a statement's solved result and may be narrowed by
selectors: \texttt{x.edge}, \texttt{x.inner}, \texttt{x.side(toward: y)},
\texttt{x.front}, \texttt{x.door} (a door anchor recorded on the art), or
\texttt{x[*]} for all instances. A statement whose region is \texttt{p[*]} is replicated
per instance, so \texttt{place p[*].house in p[*]} gives every parcel its own house.
Region expressions combine cell sets by union, intersection and difference.

\begin{table}[t]
\caption{PLL statements and their solving step}
\label{tab:pll}
\centering
\footnotesize
\begin{tabular}{@{}p{1.25cm} >{\ttfamily}p{2.85cm} p{3.5cm}@{}}
\hline
\textbf{Kind} & \textrm{\textbf{Statements}} & \textbf{Meaning (step)} \\
\hline
givens & fixed, anchor & existing or locked content; named points (a) \\
networks & path, network & least-cost roads and rivers; street networks (b) \\
areas & partition, subdivide, floorplan, graph & zones by interest maps; parcels; rooms; mission graphs (c) \\
objects & place (primary) & a region's main object (d) \\
areas & region, connect & nested areas; doors and openings (e) \\
objects & place & other objects by relations (f) \\
linear & line & fences, walls, hedges on edges (g) \\
patterns & rows, grid & beds, crops, pews; orchards (h) \\
dressing & scatter & density-driven Poisson-disk scatter (i) \\
ground & fill & autotiled terrain (j) \\
rules & relation, objective, check & constraints; soft costs; post-solve checks (k) \\
\hline
\end{tabular}
\end{table}

\paragraph{Relations} Relations are n-ary, as in SceneCraft \cite{scenecraft}, and each
has a score $s \in [0,1]$ computed from solved results. The vocabulary covers
containment (\texttt{inside}, \texttt{outside}), global position (\texttt{center},
\texttt{edge}, \texttt{corner}), distance (\texttt{adjacent}, \texttt{near},
\texttt{far}), relative position in the facing frame of an object
(\texttt{in-front-of}, \texttt{behind}, \texttt{beside}, \texttt{between}), facing
(\texttt{faces}), groups (\texttt{aligned}, \texttt{in-row}, \texttt{symmetric},
\texttt{repeat}), linear features and networks (\texttt{along}, \texttt{encloses},
\texttt{crosses}), reachability and visibility (\texttt{path-between},
\texttt{visible-from}, \texttt{focus}), stacking (\texttt{on}) and interest maps
(\texttt{interest}). For example, with $d$ the distance between members,
\begin{equation}
s_{\text{near}}(d; d_0, d_1) =
\begin{cases} 1 & d \le d_0 \\ \dfrac{d_1 - d}{d_1 - d_0} & d_0 < d < d_1 \\ 0 & d \ge d_1, \end{cases}
\end{equation}
and \texttt{interest} maps a criterion value in $[-1,1]$ to $[0,1]$, with $-1$ a hard
veto \cite{emilien2012}. Containment, enclosure, reachability, stacking and vetoes are
hard by default; any relation can be made hard. New relation types are registered in the
ontology with their arity, arguments and score function, so adding one does not change
the parser.

\paragraph{Solving order} A fixed order is necessary because some relations refer to
objects that must exist first: a garden region lies \emph{behind} the house of its plot,
so the house must be placed before the garden is solved. The solver therefore works
top-down through the region tree, and inside each region it takes statements in the
fixed steps of Table~\ref{tab:pll}: givens; paths (rivers, then main roads, streets,
lanes and trails); the region's areas; its primary objects, large before small; nested
regions and openings, which can now relate to the primary objects; other objects by a
greedy, anchor-first and large-first placement \cite{holodeck} followed by simulated
annealing, which outperformed greedy and genetic search for GLDL \cite{gldl}; lines;
rows and grids; and scatter on the remaining free cells. It then recurses into each area
produced, largest first. Terrain fill and checks run once the tree is complete. A
reference to a result solved later is a static error. Each statement's random stream is
seeded with a hash of the document seed and its id, so adding or removing one statement
does not change the others.

\paragraph{Objective} Within a region, the soft objective combines the relation scores
with GLDL's cluster and sparsity terms \cite{gldl}, a variety term and a coverage term:
\begin{equation}
\begin{aligned}
\text{cost} = &\textstyle\sum_{r \in \text{soft}} w_r (1 - s_r)
  + w_c \sum_{i<j} \frac{1}{d_{ij} + 1}
  + w_s \max_{c \in F} d(c) \\
 &+ w_v V + w_k \,|\rho - \rho^{*}|,
\end{aligned}
\end{equation}
where $F$ is the set of free cells and $d(c)$ the distance from a free cell to the
nearest object, so large empty voids are penalised; $V$ is the share of neighbouring
instances that use the same asset; and $\rho$ is the region's footprint density against
its target band $\rho^{*}$, as in CityGenAgent \cite{citygenagent}. Default weights come
from the module via the ontology.

\paragraph{Hard rules and checks} For every statement: results lie inside their region,
no two footprints overlap on a layer (fixed cells included), objects stand on terrain
their asset allows, cells marked \texttt{keepFree} (the front of a door, a stall or a
chest) stay walkable, every entrance, gate and door is reachable from the main entrance,
locked openings open only once their key is reachable, lines are continuous with valid
joins, and every asset is perspective-compatible and has an allowed licence. Because GLDL
found unplayable levels despite its constraints \cite{gldl}, reachability is tested by
pathfinding, never assumed. \texttt{check} statements add post-solve predicates
(\texttt{count}, \texttt{density}, \texttt{area}, \texttt{footprint-density},
\texttt{reachable}, \texttt{closed}, \texttt{score}, \texttt{filled}). Each error carries
a code, the statement id and its Spec node, which routes repair to the right place.

\paragraph{Example} Fig.~\ref{fig:garden} shows the text form of one garden compiled
from the garden module's template. The canonical form is JSON; the text form maps to it
one-to-one and is used in review and logs.

\begin{figure}[t]
\scriptsize
\begin{verbatim}
region   g14 in p14 shape rect size 4x3   @plot#14/garden
relation behind(g14, p14.house) weight 3
relation adjacent(g14, p14.house.back) weight 2
anchor   g14.gate at g14.side(toward: p14.house.back)
         kind gate
line     g14.fence slot fence along g14.edge
         gaps [g14.gate] essential false
rows     g14.beds in g14.inner slot bed count 2
scatter  g14.flowers in g14.inner - g14.beds slot plant
         mix {rose:.5, tulip:.3, daisy:.2} density 0.3
relation path-between(g14.gate, p14.house.back) hard
check    count(g14.beds, [1, 3])
check    closed(g14.fence)
\end{verbatim}
\vspace{-6pt}
\caption{A garden behind the house of plot 14, in PLL text form. The density 0.3 is the
numeric value of the Spec's \emph{common}.}
\label{fig:garden}
\end{figure}

\paragraph{Relation to GLDL} PLL keeps GLDL's core terms with the same meaning:
\emph{PredefinedFacility} is \pll{fixed}, \emph{NewFacility} is \pll{place} (with
\pll{region} for areas), \emph{Points} are \pll{anchor}s, its basic and relationship
constraints map to \pll{near}, \pll{inside}, \pll{edge}, \pll{far},
\pll{visible-from}, \pll{focus}, \pll{aligned} and \pll{faces}, and its objective is
the cluster and sparsity part of ours. GLDL, however, describes one 3D room of cuboid
facilities, without a hierarchy of parts, quantities or qualities, linear features,
networks, subdivision, scatter, terrain, tiles, perspectives, group relations or
provenance. PLL adds these, together with rooms, mission graphs and openings with
lock--key reachability, and in-document checks. A single room without nesting can be
exported to GLDL for comparison.

\subsection{Asset Management}
\label{sec:assets}

\paragraph{Understanding the art} Every asset in the project is analysed into a catalog
entry. Pixel analysis of the alpha mask gives the \emph{footprint} in cells (the opaque
base, not the image size) and the \emph{anchor} (the lowest opaque row and its shadow);
edge analysis groups fence, wall and hedge sprites into \emph{piece sets} (straight,
corner, end, T, cross, gate); a vision model, run on an upscaled image as recommended by
GameTileNet \cite{gametilenet}, gives the category, tags, the anchors on the art (door,
gate, seat) and the \emph{perspective} (top-down, three-quarter or isometric). Each entry
also records the facings the art supports, its palette, layer, passability, its
affordance and tile type from GameTileNet's taxonomy, and its licence, author and source.
Tilesheets are segmented into whole objects before analysis. Flat top-down art and
upright art may be mixed in one project, as most packs do; what must match is the grid,
since top-down linear pieces join only on a square grid and isometric pieces only on an
isometric one.

\paragraph{Needs} From the confirmed Brief, the asset manager walks the slots of every
enabled component and its descendants, resolves each slot's role and concepts from the
component's quality fields (a flower garden with a rose and tulip mix needs
\texttt{flower:rose} and \texttt{flower:tulip}), and merges identical needs across
components. Each need gets an importance: \emph{critical} (essential and large or
structural: buildings, terrain, roads, walls, fences), \emph{normal} (essential, small)
or \emph{optional} (adaptable). Buildings, terrain, piece sets and things seen many times
are \emph{style-critical}.

\paragraph{Fit} A candidate is any asset from a source that passes the hard filters: it
fills the role and shows the concept (any rose fills \emph{flower} through
\texttt{broader}); its perspective is compatible with the grid; its footprint fits the
need's range; its layer is right; its licence is known and allowed; and, for a piece
set, the pieces a line needs are present or can be generated. Because text similarity
alone matched affordances poorly in \cite{chen2025narrative}, these filters come before
any similarity. The remaining candidates are scored
\begin{equation}
\text{fit} = 0.40\,\sigma_{\text{sem}} + 0.35\,\sigma_{\text{style}} + 0.15\,\sigma_{\text{size}} + 0.10\,\sigma_{\text{qual}},
\end{equation}
where $\sigma_{\text{sem}}$ measures how well the asset shows the concept,
$\sigma_{\text{style}}$ its closeness to the place's \emph{style profile} (palette
distance, pixels per cell, outline weight and shading, measured from the assets the place
already uses), $\sigma_{\text{size}}$ how little rescaling it needs and
$\sigma_{\text{qual}}$ its cleanliness and completeness. For style-critical needs,
$\sigma_{\text{style}} \ge 0.7$ is also required. A fit of at least 0.75 is \emph{good};
from 0.55 it is \emph{usable} after adaptation (recolouring to the palette, rescaling,
mirroring); below that it is rejected.

\paragraph{Sources and decision procedure} Four sources share one interface: the
\emph{project}; a curated, pre-analysed \emph{library} organised in style families;
the \emph{web}, restricted to approved free game-art sites; and \emph{generation} with
the editor's image tool. Under the default \emph{auto} policy, needs are resolved in
order of importance as follows: use a good project candidate; otherwise import a good
library candidate; otherwise, for a style-critical need, generate it, since buildings,
terrain and piece sets fetched from the web rarely match a pack; otherwise search the web
if the project allows it, and else generate. If nothing good was found, the best usable
candidate is adapted and reported; failing that, an optional need is dropped and an
essential one receives a labelled placeholder, so the part still builds and the report
lists it. The user can narrow the order per component (project only, library, web or
generate) in the second confirmation round, where the asset plan is shown as a table of
needs with their proposed sources, thumbnails and alternatives.

Generation prompts the image model with the need (role, concept, the module's and
Brief's descriptions, footprint, perspective and layer), the style profile and one to
three reference images from the project closest in role. The result is analysed like any
asset and judged by a critic for concept, style, footprint and perspective, with up to
three attempts, each including the critic's reasons. For partial piece sets, the missing
pieces are generated from an existing piece as reference, and their joins are verified
by edge analysis. Web retrieval builds queries from the concept, its ontology synonyms,
the perspective and style words; downloads candidates into an isolated folder; reads the
licence and author from the asset's page; segments and analyses them; and scores them as
above. A candidate whose licence cannot be read is rejected. Allowed licences are a
project setting (by default CC0, CC-BY and OGA-BY), assets that need attribution are
added to the project's credits automatically, and verification checks that every built
entity's asset has an allowed licence. All needs are resolved before compiling, and
independent generation and web requests run in parallel.

\section{Evaluation}
% To be written: ablation of components (ontology subgraph, confirmation, expansion
% rules, verification and repair, asset procedure) on the same queries, with
% human ranking of the generated places.

% \section*{Acknowledgment}

\begin{thebibliography}{00}
\bibitem{parish2001} Y. I. H. Parish and P. M\"uller, ``Procedural modeling of cities,'' in \textit{Proc. SIGGRAPH}, 2001, pp. 301--308.
\bibitem{shaker2016pcg} N. Shaker, J. Togelius, and M. J. Nelson, \textit{Procedural Content Generation in Games}. Cham, Switzerland: Springer, 2016.
\bibitem{todd2023} G. Todd, S. Earle, M. U. Nasir, M. C. Green, and J. Togelius, ``Level generation through large language models,'' in \textit{Proc. Int. Conf. Foundations of Digital Games (FDG)}, 2023.
\bibitem{sudhakaran2023} S. Sudhakaran, M. Gonz\'alez-Duque, M. Freiberger, C. Glanois, E. Najarro, and S. Risi, ``MarioGPT: Open-ended text2level generation through large language models,'' in \textit{Proc. NeurIPS}, 2023.
\bibitem{holodeck} Y. Yang \textit{et al.}, ``Holodeck: Language guided generation of 3D embodied AI environments,'' in \textit{Proc. IEEE/CVF Conf. Computer Vision and Pattern Recognition (CVPR)}, 2024.
\bibitem{scenecraft} Z. Hu, A. Iscen, A. Jain, T. Kipf, Y. Yue, D. A. Ross, C. Schmid, and A. Fathi, ``SceneCraft: An LLM agent for synthesizing 3D scenes as Blender code,'' in \textit{Proc. Int. Conf. Machine Learning (ICML)}, 2024.
\bibitem{gldl} K. Xu and C. Verbrugge, ``Constraint is all you need: Optimization-based 3D level generation with LLMs,'' in \textit{Proc. Int. Conf. Foundations of Digital Games (FDG)}, 2025.
\bibitem{togelius2011} J. Togelius, G. N. Yannakakis, K. O. Stanley, and C. Browne, ``Search-based procedural content generation: A taxonomy and survey,'' \textit{IEEE Trans. Comput. Intell. AI Games}, vol. 3, no. 3, pp. 172--186, 2011.
\bibitem{emilien2012} A. Emilien, A. Bernhardt, A. Peytavie, M.-P. Cani, and E. Galin, ``Procedural generation of villages on arbitrary terrains,'' \textit{The Visual Computer}, vol. 28, pp. 809--818, 2012.
\bibitem{vanegas2012} C. A. Vanegas, T. Kelly, B. Weber, J. Halatsch, D. G. Aliaga, and P. M\"uller, ``Procedural generation of parcels in urban modeling,'' \textit{Computer Graphics Forum}, vol. 31, no. 2, pp. 681--690, 2012.
\bibitem{karth2017} I. Karth and A. M. Smith, ``WaveFunctionCollapse is constraint solving in the wild,'' in \textit{Proc. Int. Conf. Foundations of Digital Games (FDG)}, 2017.
\bibitem{alaka2023} S. Alaka and R. Bidarra, ``Hierarchical semantic wave function collapse,'' in \textit{Proc. Int. Conf. Foundations of Digital Games (FDG)}, 2023.
\bibitem{dormans2010} J. Dormans, ``Adventures in level design: Generating missions and spaces for action adventure games,'' in \textit{Proc. Workshop on Procedural Content Generation in Games (PCG)}, 2010.
\bibitem{layoutgpt} W. Feng \textit{et al.}, ``LayoutGPT: Compositional visual planning and generation with large language models,'' in \textit{Proc. NeurIPS}, 2023.
\bibitem{infinigen} A. Raistrick \textit{et al.}, ``Infinigen Indoors: Photorealistic indoor scenes using procedural generation,'' in \textit{Proc. IEEE/CVF Conf. Computer Vision and Pattern Recognition (CVPR)}, 2024.
\bibitem{3dgpt} C. Sun, J. Han, W. Deng, X. Wang, Z. Qin, and S. Gould, ``3D-GPT: Procedural 3D modeling with large language models,'' in \textit{Proc. Int. Conf. 3D Vision (3DV)}, 2025.
\bibitem{worldcraft} X. Liu, C.-K. Tang, and Y.-W. Tai, ``WorldCraft: Photo-realistic 3D world creation and customization via LLM agents,'' 2025, arXiv:2502.15601.
\bibitem{citygenagent} Z. Liu \textit{et al.}, ``Imagine a city: CityGenAgent for procedural 3D city generation,'' 2026, arXiv:2602.05362.
\bibitem{xu2025database} K. Xu and C. Verbrugge, ``A database-driven framework for 3D level generation with LLMs,'' in \textit{Proc. AAAI Conf. Artificial Intelligence and Interactive Digital Entertainment (AIIDE)}, vol. 21, 2025, pp. 367--376.
\bibitem{chen2025narrative} Y.-C. Chen and A. Jhala, ``Narrative-to-scene generation: An LLM-driven pipeline for 2D game environments,'' 2025, arXiv:2509.04481.
\bibitem{gametilenet} Y.-C. Chen and A. Jhala, ``GameTileNet: A semantic dataset for low-resolution game art in procedural content generation,'' in \textit{Proc. AAAI Conf. Artificial Intelligence and Interactive Digital Entertainment (AIIDE)}, vol. 21, 2025, pp. 12--21.
\bibitem{tutenel2009} T. Tutenel, R. M. Smelik, R. Bidarra, and K. J. de Kraker, ``Using semantics to improve the design of game worlds,'' in \textit{Proc. AAAI Conf. Artificial Intelligence and Interactive Digital Entertainment (AIIDE)}, 2009.
\bibitem{smelik2011} R. M. Smelik, T. Tutenel, K. J. de Kraker, and R. Bidarra, ``A declarative approach to procedural modeling of virtual worlds,'' \textit{Computers \& Graphics}, vol. 35, no. 2, pp. 352--363, 2011.
\bibitem{owl2} W3C OWL Working Group, ``OWL 2 Web Ontology Language document overview (second edition),'' W3C Recommendation, Dec. 2012.
\bibitem{skos} A. Miles and S. Bechhofer, Eds., ``SKOS Simple Knowledge Organization System reference,'' W3C Recommendation, Aug. 2009.
\bibitem{shacl} H. Knublauch and D. Kontokostas, Eds., ``Shapes Constraint Language (SHACL),'' W3C Recommendation, Jul. 2017.
\bibitem{clareval} J. Li, Y. Wu, and Y. Chang, ``ClarEval: A benchmark for evaluating clarification skills of code agents under ambiguous instructions,'' 2026, arXiv:2603.00187.
\bibitem{kobalczyk2025} K. Kobalczyk, N. Astorga, T. Liu, and M. van der Schaar, ``Active task disambiguation with LLMs,'' in \textit{Proc. Int. Conf. Learning Representations (ICLR)}, 2025.
\bibitem{tam2024} Z. R. Tam \textit{et al.}, ``Let me speak freely? A study on the impact of format restrictions on performance of large language models,'' in \textit{Proc. EMNLP: Industry Track}, 2024.
\bibitem{watabou} Watabou, ``Medieval fantasy city generator,'' 2016. [Online]. Available: https://watabou.itch.io/medieval-fantasy-city-generator
\bibitem{lynch1960} K. Lynch, \textit{The Image of the City}. Cambridge, MA: MIT Press, 1960.
\end{thebibliography}

\end{document}
